% ============================================================
% Chapter 4: Results
% ============================================================
\chapter{Results}
\label{ch:results}

\section{Overview}
\label{sec:results-overview}

This chapter presents the findings of the research. Results are organized
by research question or experiment, and are presented objectively without
interpretation (reserved for the Discussion chapter).

\section{Experiment 1: \emph{Description}}
\label{sec:exp1}

\subsection{Setup}
\label{sec:exp1-setup}

Describe the specific configuration for this experiment.

\subsection{Results}
\label{sec:exp1-results}

Present the results using tables and figures.

\begin{table}[htbp]
    \centering
    \caption{Results for Experiment 1}
    \label{tab:exp1}
    \begin{tabular}{lccc}
        \toprule
        \textbf{Method} & \textbf{Metric 1} & \textbf{Metric 2} & \textbf{Metric 3} \\
        \midrule
        Baseline        & 0.00 & 0.00 & 0.00 \\
        Proposed        & 0.00 & 0.00 & 0.00 \\
        \bottomrule
    \end{tabular}
\end{table}

\begin{figure}[htbp]
    \centering
    % \includegraphics[width=0.8\textwidth]{example-figure}
    \caption{Visualization of Experiment 1 results}
    \label{fig:exp1}
\end{figure}

\section{Experiment 2: \emph{Description}}
\label{sec:exp2}

\subsection{Setup}
\label{sec:exp2-setup}

Describe the specific configuration for this experiment.

\subsection{Results}
\label{sec:exp2-results}

Present the results using tables and figures.

\section{Summary of Results}
\label{sec:results-summary}

Provide a concise summary of all results, highlighting patterns and
trends observed across experiments.
